\documentclass[11pt]{article}
\usepackage[T1]{fontenc}
\usepackage{textcomp}
\usepackage[margin=0.75in]{geometry}
\usepackage{amsmath,amssymb,amsthm}
\usepackage{array,booktabs}
\usepackage{hyperref}
\setlength{\parindent}{0pt}
\newtheorem{theorem}{Theorem}
\theoremstyle{definition}
\newtheorem{definition}{Definition}
\title{Flow Proof Artifact: inscribed-on-diameter-right-alt}
\author{Generated by \texttt{flow doc proof}}
\date{Flow Proof Book}
\begin{document}
\maketitle
An inscribed angle on a diameter is right, alternate witness.\\[0.5em]
\textbf{Source.} euclid — Elements, Book III, Proposition 31\\[0.5em]
\bigskip
\noindent{\large \textbf{Derived fact 1.} \textit{angle APB equals one right angle on diameter AB} \hfill the Euclidean plane \cdot circle \cdot inscribed on diameter right alt}\par
\smallskip\noindent\textit{Coordinate: the Euclidean plane \cdot circle \cdot inscribed on diameter right alt}\par
\smallskip\noindent\textit{Source: euclid}\par
\smallskip\noindent\textit{Built on: inscribed on semicircle is right, for circles in on the Euclidean plane, inscribed angle on diameter is right, for circles in on the Euclidean plane}\par
\medskip
\noindent\textbf{Goal.} angle APB equals one right angle on diameter AB\par
\noindent$\displaystyle \angle APB = 90^\circ$\par
\medskip
\noindent
\begin{tabular}{@{} >{\raggedright\arraybackslash}p{0.06\textwidth} >{\raggedright\arraybackslash}p{0.47\textwidth} >{\raggedleft\arraybackslash}p{0.41\textwidth} @{}}
\textbf{\#} & \textbf{Proof} & \textbf{Mathematics} \\
\midrule
\textbf{1.} & We prove that inscribed on diameter right alt for circles in the Euclidean plane. & \\[0.45em]
\textbf{2.} & We invoke the derived fact governing circles in the Euclidean plane: inscribed on semicircle is right, for circles in on the Euclidean plane. & \\[0.45em]
\textbf{3.} & We invoke the derived fact governing circles in the Euclidean plane: inscribed angle on diameter is right, for circles in on the Euclidean plane. & \\[0.45em]
\textbf{4.} & From step 2 and step 3, this implies angle APB equals one right angle. Hence proven. & $\displaystyle \angle APB = 90^\circ$ \\[0.45em]
\bottomrule
\end{tabular}
\label{thm:Geometry-circle-inscribed_on_diameter_right_alt}
\par\medskip\hrule
\end{document}
